\section{Introduction}
\label{sec:introduction}

% TARGET: ~1.25 pages
% Structure: (1) Motivation/problem, (2) Existing approaches and their limitations,
%            (3) Our approach, (4) Contributions (bulleted)

% PLACEHOLDER — to be written

% Paragraph 1: Clinical ML models fail across hospital systems (domain shift)
% - EHR data varies across institutions (equipment, protocols, coding)
% - Models trained on one system degrade on another
% - Cite: YAIB benchmark showing cross-domain degradation

% Paragraph 2: Existing DA approaches and their limitations
% - Feature-space methods (DANN, CORAL) modify internal representations
% - Require retraining or access to model internals
% - Don't preserve clinical interpretability of input features

% Paragraph 3: Our approach — input-space translation
% - Key insight: translate INPUT data, not features
% - Frozen target model = practical deployment scenario
% - Retrieval-guided: leverage target domain examples for instance-level adaptation

% Paragraph 4: Contributions (bulleted)
\paragraph{Contributions.}
\begin{itemize}
    \item We formulate \textbf{input-space domain adaptation} for clinical time series, where a frozen target-trained model is adapted by transforming source inputs rather than modifying model internals.
    \item We propose \retr{}, a retrieval-guided translation architecture that uses a shared encoder, target-domain memory bank, and cross-attention to produce instance-level adapted inputs.
    \item We demonstrate state-of-the-art results across \textbf{5 clinical tasks} (mortality, AKI, sepsis, length of stay, kidney function), surpassing 8 DA baselines and matching or exceeding native-domain performance.
    \item We show the translator is \textbf{architecture-agnostic}: translations trained for an LSTM transfer to GRU and TCN models without retraining, preserving 47--86\% of the improvement.
    \item We validate \textbf{multi-source generalization}: a second source domain (HiRID$\to$MIMIC) yields even larger gains, confirming the method is not source-specific.
\end{itemize}
